\begin{textsample}{POS Dim 1 – vem\_ed\_31 – Score 18.00 – vem\_ed\_31/t163.txt}  \label{ex:f1_pos_040}
Preparando alunos para profissões \textbf{que} não existem “Prepping COIL students \textbf{for} jobs not yet created”.
Os projetos COIL (Collaborative Online International Learning, ou Aprendizagem Colaborativa Internacional On-line) contribuindo para preparar estudantes para empregos \textbf{que} ainda não \textbf{foram} \textbf{criados}.
Esse \textbf{foi} o tema do SUNY COIL Community Webinar conduzido em 17 de setembro de 2025 por Osvaldo Succi Junior, coordenador da equipe de Apoio à Internacionalização do Ensino Superior da Depes / CGESG / CPS, e Luciane Stallivieri, da Universidade Federal de Santa Catarina.
Luciane Stallivieri \textbf{propõe} a noção de Internacionalização Responsável, \textbf{que} busca equilíbrio (\textbf{linguístico}, geográfico, econômico) na escolha de parceiros, compartilha resultados com a sociedade (accountability), oferece inclusão e garante sustentabilidade aos \textbf{processos}, com comprometimento \textbf{institucional} (compliance).
Para encarar o “VUCA world” (VUCA \textbf{é} um acrônimo em inglês para volatilidade, incerteza, complexidade e ambiguidade) \textbf{é} preciso ter visão / \textbf{foco} no propósito, \textbf{compreensão} para refinar / adaptar planos, clareza na criação desses planos e agilidade / \textbf{flexibilidade} e colaboração para alcançar os objetivos.
Considerando as macrotendências das \textbf{transformações} no mundo dos negócios, as incertezas, as \textbf{mudanças} aceleradas e os empregos a \textbf{serem} \textbf{criados} (e destruídos), Osvaldo Succi \textbf{defende} \textbf{que} a educação superior deve evoluir para preparar os estudantes com \textbf{habilidades} \textbf{que} transcendam as descrições \textbf{específicas} de cargos, \textbf{promovendo} adaptabilidade, pensamento \textbf{crítico} e envolvimento ético global.
Afinal, as principais \textbf{competências} centrais no mundo atual \textbf{são}: pensamento analítico; resiliência, \textbf{flexibilidade} e agilidade; liderança; pensamento criativo; motivação e autoconsciência.
Após um exercício criativo com profissões \textbf{que} ainda não existem, como “corretor-guardião de dados pessoais” ou “terapeuta de robôs”, Osvaldo Succi listou algumas \textbf{competências} desejáveis a \textbf{serem} incorporadas no design de projetos COIL (Collaborative Online International Learning, ou Aprendizagem Colaborativa Internacional On-line): Sensibilidade lógica humano-computador Resumo confiável de \textbf{conteúdo} Design ético e auditoria de IA Avaliação de riscos \textbf{Desenvolvimento} de cenários plausíveis Escalabilidade contextual Aprendizagem e desaprendizagem adaptativas Pensamento sistêmico Gestão da privacidade

% matched lemmas: competência, compreensão, conteúdo, criar, crítico, defender, desenvolvimento, específico, flexibilidade, foco, habilidade, institucional, linguístico, mudança, processo, promover, propor, que, ser, transformação
\end{textsample}
